% Synthetic fixture: an ACM conference paper.
\documentclass[sigconf,nonacm]{acmart}
\settopmatter{printacmref=false}
\renewcommand\footnotetextcopyrightpermission[1]{}
\pagestyle{plain}

\begin{document}

\title{Reading Receipts for Shared Notebooks}

\author{Ines Carvalho}
\affiliation{%
  \institution{Example Institute}
  \city{Porto}
  \country{Portugal}}
\email{ines@example.org}

\author{Tomasz Wrona}
\affiliation{%
  \institution{Example Institute}
  \city{Gdansk}
  \country{Poland}}
\email{tomasz@example.org}

\begin{abstract}
Shared notebooks lose track of who has seen which change. We add a small
receipt to every cell and show that readers notice edits sooner.
\end{abstract}

\begin{CCSXML}
<ccs2012>
 <concept>
  <concept_id>10003120.10003130</concept_id>
  <concept_desc>Human-centered computing~Collaborative and social computing</concept_desc>
  <concept_significance>500</concept_significance>
 </concept>
</ccs2012>
\end{CCSXML}

\ccsdesc[500]{Human-centered computing~Collaborative and social computing}

\keywords{notebooks, awareness, collaboration}

\begin{teaserfigure}
  \centering
  \fbox{\parbox{0.8\textwidth}{\centering A notebook cell with its receipt}}
  \caption{Each cell shows who has read its latest version.}
  \label{fig:teaser}
\end{teaserfigure}

\maketitle

\section{Introduction}
Teams share notebooks for analysis, and a change made late in the evening is
often missed \cite{prieto20}. Figure~\ref{fig:teaser} shows our remedy.

\section{Design}
A receipt records the reader and the version. We keep it under
40~bytes\footnote{Measured on the final prototype.} per cell, and show it
as a small dot.

\begin{itemize}
  \item A grey dot means nobody has read the cell.
  \item A filled dot means every collaborator has.
\end{itemize}

\section{Study}
Twelve people used the notebooks for two weeks. Edits were noticed after a median
of 9 minutes instead of 40.

\bibliographystyle{ACM-Reference-Format}
\bibliography{refs}

\end{document}
